\section{Исследование рынка}

Исследование сегмента рынка, на который ориентирован разрабатываемый
продукт, направлено на определение области его применения, оценку
потенциального спроса и прогнозирование количества возможных продаж.

Разрабатываемый продукт представляет собой специализированный
инструмент и методику верификации соответствия реализации формальным
спецификациям на языке TLA+. Основной областью применения являются
организации, разрабатывающие сложные программные системы с высокими
требованиями к корректности: распределённые системы, финансовые
сервисы, системы хранения данных и инфраструктурные решения.

В настоящее время применение формальных методов верификации остаётся
ограниченным из-за высокой сложности внедрения, необходимости
квалифицированных специалистов и отсутствия удобных инструментов
интеграции с существующим кодом. В частности, использование TLA+
требует:
\begin{itemize}
\item разработки формальной спецификации системы;
\item сопоставления абстрактной модели с реальной реализацией;
\item инструментирования кода для получения трасс исполнения;
\item анализа результатов модельной проверки.
\end{itemize}

Разрабатываемый продукт решает задачу сопоставления трасс исполнения с
формальной спецификацией, однако его внедрение остаётся сложным
процессом, требующим участия квалифицированных разработчиков и
аналитиков. Это существенно ограничивает массовость применения и
снижает число потенциальных покупателей, одновременно увеличивая
ценность продукта для каждой конкретной организации.

По оценкам, ежегодно реализуется порядка 3000 проектов, из которых
около 5\% могут потенциально использовать формальные методы
верификации:

\begin{equation}
N_p = 0.05 \cdot 3000 = 150 \text{ проектов в год}.
\end{equation}

С учётом высокой сложности внедрения, специфичности продукта и
необходимости адаптации под конкретные системы, примем долю реальных
покупок на уровне 10\% от потенциального спроса:

\begin{equation}
N_p^o = 0.1 \cdot 150 = 15 \text{ продаж в год}.
\end{equation}

Средняя стоимость одной копии продукта определяется его
специализированным характером и включает не только стоимость
программного обеспечения, но и затраты на внедрение, интеграцию и
сопровождение. В отличие от массовых программных продуктов, подобные
решения поставляются как комплексный инструмент, ориентированный на
решение критически важных задач.

Высокая стоимость обусловлена следующими факторами:
\begin{itemize}
\item необходимость адаптации инструмента под конкретную систему;
\item сложность разработки и сопровождения TLA+ спецификаций;
\item высокая квалификация специалистов, участвующих во внедрении;
\item значительная ценность продукта для предотвращения ошибок в
критических системах.
\end{itemize}

С учётом указанных факторов примем среднюю стоимость одной копии:

\begin{equation}
K_{\text{пр}} = 250000 \text{ руб.}
\end{equation}

Тогда прогнозируемый объём продаж в год составит:

\begin{equation}
V = N_p^o \cdot K_{\text{пр}} =
15 \cdot 250000 = 3750000 \text{ руб.}
\end{equation}

За период в два года суммарный объём продаж составит:

\begin{equation}
V_{2} = 2 \cdot 3750000 = 750000 \text{ руб.}
\end{equation}

Таким образом, несмотря на сравнительно небольшое количество продаж,
высокая стоимость продукта обеспечивает значительный объём выручки.
Это объясняется тем, что продукт ориентирован на узкий сегмент рынка и
представляет собой высокоспециализированное решение с высокой
добавленной стоимостью.
